\chapter{Introduction}
\label{ch:Introduction}


During the last decade, aviation has become an important part of the lives of people all over the world.
Aircraft have been used for long-distance transportation between airports, connecting cities from all over the world.
However, an issue persists.
After arriving at an airport, another form of transport must be used to cover the so-called \enquote{last mile}.
Moreover, having to use airport infrastructure adds complexity to aircraft operations.
Thus, there is a pressing question that has some promising implications.
What if it were possible to remove airports from the equation?
This is where the concept of Inter-Urban Air Mobility comes in.
Inter-UAM focuses on flying inside big metropolitan areas or even between smaller cities.
The main aspects are a longer range than Urban Air Mobility,\qty{>100}{km} versus $\approx$\qty{30}{km}, and the ability to take off and land vertically from vertiports in urban areas.

 This report aims to choose a promising inter-UAM design focusing on reducing both the ground- and environmental footprint of the aircraft.
 This means that the vehicle could take off from smaller vertiports or even driveways, making \enquote{last-mile delivery} possible.
 Additionally, the vehicle will be fully electric to facilitate a greener future.
 This leads to the following Mission Need Statement, which will be achieved through the following Project Object Statement:

\paragraph{MNS:}“\textit{Achieve Sustainable Inter-Urban Air Mobility with a low ground footprint vehicle.}”
\vspace{-5mm}
\paragraph{POS:}“\textit{Design a low ground footprint, sustainable, urban transwing electric Vertical Take-Off and Landing vehicle within production costs of 2000000 EUR, by 10 students in 10 weeks time.}”

To ensure that the project runs smoothly and is finished on time, the Organisation of the team and the Logic diagrams are discussed in \Cref{ch:Organisation}.
Afterwards, the mission profile is sketched in \Cref{ch:mission_profile}, from which different design options are explored in \Cref{ch:Concepts}.
To enable proper comparison between the different concepts, they are preliminarily sized in \Cref{ch:PrelimSizing}.
After which, the trade-off procedure is illustrated in \Cref{ch:TradeOff}.
This leads to a final design choice that will be pursued in the near future.
However, it is important to ensure that the design choice is correct and the trade-off is not too sensitive.
This is why a sensitivity analysis is conducted in \Cref{ch:SensitivityAnalysis}, followed by verification and validation of the tools used for the preliminary design, which can be found in \Cref{ch:VerificationValidation}.
Another important aspect that shall be tackled in this report is the interfaces between the constituting elements of the design, this is detailed in \Cref{ch:Interfaces}.
Moreover, to ensure the success of the project, the most important risks are assessed and defined in \Cref{chapter:TechRiskAnalysis}.
Lastly, it is important to ensure that the final product is up to standard and, therefore, sustainable.
This is done by developing a Production Plan in \Cref{ch:ProductionPlan}, while the strategy to design a sustainable product is detailed in \Cref{ch:sustainability}.